% part: intuitionistic-logic
% chapter: propositions-as-types
% section: normalization

\documentclass[../../../include/open-logic-section]{subfiles}

\begin{document}

\olfileid{int}{pty}{nor}

\olsection{સામાન્યીકરણ}

આ વિભાગમાં આપણે સાબિત કરીશું કે ઘટાડાના
કોઈ ક્રમથી દરેક નિગમનને સામાન્ય નિગમનમાં
ઘટાડી શકાય છે. તેને \emph{સામાન્યીકરણનો ગુણધર્મ}
કહે છે. આપણે પ્રતિજ્ઞાપનો અને પ્રકારો વચ્ચેની
અનુરૂપતાનો ઉપયોગ કરીશું: દરેક સાબિતી-પદને
સામાન્ય સ્વરૂપમાં ઘટાડી શકાય છે તે બતાવીશું;
પછી સ્વાભાવિક નિગમનની !!{derivation}sનું
સામાન્યીકરણ અનુસરે છે.

પહેલા પદોની જટિલતા માપતાં કેટલાંક વિધેયો
વ્યાખ્યાયિત કરીએ. !!a{formula}ની
\emph{લંબાઈ}~$\len{!A}$ આ પ્રમાણે છે:
\begin{align*}
  \len{p} &= 0\\
  \len{!A \land !B} &= \len{!A} + \len{!B} + 1\\
  \len{!A \lor !B} &= \len{!A} + \len{!B} + 1 \\
  \len{!A \lif !B} &= \len{!A} + \len{!B} + 1.
\end{align*}
સંકોચ્ય પદ~$M$ની જટિલતા તેના
\emph{કટ-ક્રમ}~$\cutrank{M}$થી માપીએ:
\begin{align*}
  \cutrank{(\lambd[\typeof{x}{!A}][\typeof{N}{!B}])Q} &=
  \len{!A} + \len{!B} + 1 \\
  \cutrank{\ande{i}{\andi{\typeof{M}{!A}}{\typeof{N}{!B}}}} &=
  \len{!A} + \len{!B} + 1 \\
  \cutrank{\ore{\ori{i}{!A_{i}}{\typeof{M}{!A_i}}}
    {\typeof{x_1}{!A_1}}{\typeof{N_1}{!C}}
    {\typeof{x_2}{!A_2}}{\typeof{N_2}{!C}}} &=
  \len{!A_1} + \len{!A_2} + 1
\end{align*}
\sourcecorrection{OLMD-280}{મૂળની સરવાળા-પ્રકારની છેલ્લી વ્યાખ્યામાં અસંબંધિત $!A$ અને $!B$ આવે છે; ઉપરના પદના પ્રકારો મુજબ $!A_1$ અને $!A_2$ કર્યા છે.}
સાબિતી-પદની જટિલતા તેમાંના સૌથી જટિલ
સંકોચ્ય ઉપપદથી માપીએ; પદ સામાન્ય હોય,
તો માપ~$0$ છે:
\[
  \maxrank{M} = \max \bigl(\{0\} \cup \{\cutrank{N} | N \text{ એ } M \text{
    નું સંકોચ્ય ઉપપદ છે}\}\bigr)
\]
\sourcecorrection{OLMD-286}{મૂળ ગદ્ય સામાન્ય પદનો મહત્તમ કટ-ક્રમ શૂન્ય કહે છે, પરંતુ દર્શાવેલું સૂત્ર ખાલી ગણનું મહત્તમ લે છે. ગદ્યના આ કિસ્સાને સૂત્રમાં પણ આવરી લેવા શૂન્ય ઉમેર્યું છે.}

\begin{lem}\ollabel{lem:subst}
  $\Subst{M}{\typeof{N}{!A}}{\typeof{x}{!A}}$ સંકોચ્ય પદ
  હોય અને $M \not\ident x$ હોય, તો આ પૈકીનો
  એક કિસ્સો બને:
  \begin{enumerate}
  \item $M$ પોતે સંકોચ્ય પદ છે, અથવા
  \item $M$નું સ્વરૂપ $\ande{i}{x}$ છે અને
    $N$નું સ્વરૂપ $\andi{P_1}{P_2}$ છે;
  \item $M$નું સ્વરૂપ $\ore{x}{x_1}{P_1}{x_2}{P_2}$ છે અને
    $N$નું સ્વરૂપ $\ori{i}{}{Q}$ છે;
    \sourcecorrection{OLMD-281}{મૂળમાં વિકલ્પ-લોપનો પ્રથમ આર્ગ્યુમેન્ટ $i$ છે; અહીં જે ચલનું પ્રતિસ્થાપન થાય છે તે $x$ હોવો જોઈએ.}
  \item $M$નું સ્વરૂપ $x Q$ છે અને
    $N$નું સ્વરૂપ $\lambd[x][P]$ છે.
  \end{enumerate}
  પહેલા કિસ્સામાં $\cutrank{\Subst{M}{N}{x}} = \cutrank{M}$;
  બીજા કિસ્સાઓમાં $\cutrank{\Subst{M}{N}{x}} = \len{!A}$ છે.
  \sourcecorrection{OLMD-282}{મૂળની છેલ્લી સમતામાં વધારાનો બંધ કૌંસ કાઢ્યો છે.}
\end{lem}
\begin{proof}
  $M$ પર આગમનથી સાબિતી કરીએ.
  \begin{enumerate}
  \item $M$ એક જ ચલ $y$ હોય અને $y \not\ident x$ હોય,
    તો $\Subst{y}{N}{x}$ એ $y$ જ છે, એટલે સંકોચ્ય નથી.
  \item $M$નું સ્વરૂપ $\andi{N_1}{N_2}$, $\lambd[x][N]$,
    અથવા $\ori{i}{!A}{N}$ હોય, તો
    $\Subst{M}{\typeof{N}{!A}}{\typeof{x}{!A}}$નું પણ
    એ જ સ્વરૂપ છે, એટલે તે સંકોચ્ય નથી.
  \item $M$નું સ્વરૂપ $\ande{i}{P}$ હોય, તો બે કિસ્સા લઈએ:
    \begin{enumerate}
      \item $P$નું સ્વરૂપ $\andi{P_1}{P_2}$ હોય, તો
        $M \ident \ande{i}{\andi{P_1}{P_2}}$ સંકોચ્ય પદ છે;
        સ્પષ્ટ રીતે
        \[
        \Subst{M}{N}{x} \ident
        \ande{i}{\andi{\Subst{P_1}{N}{x}}{\Subst{P_2}{N}{x}}}
        \]
        પણ સંકોચ્ય પદ છે. બંનેના કટ-ક્રમ સમાન છે.
      \item $P$ એક જ ચલ હોય, તો પ્રતિસ્થાપનથી સંકોચ્ય
        પદ બનાવવા તે $x$ જ હોવો જોઈએ, અને $N$નું
        સ્વરૂપ $\andi{P_1}{P_2}$ હોવું જોઈએ. હવે
        \[
        \Subst{M}{N}{x} \ident \Subst{\ande{i}{x}}{\andi{P_1}{P_2}}{x},
        \]
        લો; આ $\ande{i}{\andi{P_1}{P_2}}$ છે.
        તેનો કટ-ક્રમ પ્રકારની લંબાઈ $\len{!A}$ છે.
        \sourcecorrection{OLMD-283}{મૂળ અહીં ચલ $x$નો કટ-ક્રમ લે છે, પરંતુ કટ-ક્રમ ફક્ત સંકોચ્ય પદ માટે વ્યાખ્યાયિત છે. પ્રતિસ્થાપનથી બનેલા કટનો ક્રમ $x$ના પ્રકારની લંબાઈ છે.}
    \end{enumerate}
  \end{enumerate}
  $\ore{N}{x_1}{N_1}{x_2}{N_2}$ અને $P Q$ના કિસ્સા સમાન છે.
  \sourcecorrection{OLPG-0687-SUBST}{મૂળ સાબિતી કેટલાક કિસ્સા સમાન કહીને છોડે છે અને અધ્યાયમાં રજૂ થયેલ અસત્ય-લોપ તથા ક્રમચય-રૂપાંતરો માટે આ કટ-ક્રમની દલીલ આપતી નથી. અનુવાદ એ ખાલી જગ્યા છુપાવતો નથી; અહીંના પૂર્ણ સામાન્યીકરણ-દાવાની સ્રોત-સાબિતી અધૂરી છે.}
\end{proof}

\begin{lem}
  $M$ સંકોચાઈને $M'$ થાય અને $M$ના દરેક પોતાના કરતાં
  નાના સંકોચ્ય ઉપપદ~$N$ માટે $\cutrank{M} > \cutrank{N}$
  હોય, તો $\cutrank{M} > \maxrank{M'}$ છે.
\end{lem}

\begin{proof}
  કિસ્સાવાર સાબિતી કરીએ.
  \begin{enumerate}
  \item $M$નું સ્વરૂપ $\ande{i}{\andi{M_1}{M_2}}$ હોય,
    તો $M'$ એ $M_i$ છે. $M_i$નું દરેક ઉપપદ
    $M$નું પણ પોતાના કરતાં નાનું ઉપપદ છે,
    એટલે દાવો સાચો છે.
  \item $M$નું સ્વરૂપ
    $(\lambd[\typeof{x}{!A}][N])\typeof{Q}{!A}$ હોય, તો
    $M'$ એ $\Subst{N}{\typeof{Q}{!A}}{\typeof{x}{!A}}$ છે.
    $M'$માંનું એક સંકોચ્ય પદ લો. કાં તો~$N$માં
    તેને અનુરૂપ સંકોચ્ય પદનો કટ-ક્રમ સમાન છે,
    જે ધારણા મુજબ $\cutrank{M}$ કરતાં નાનો છે;
    અથવા તેનો કટ-ક્રમ $\len{!A}$ છે, જે વ્યાખ્યા મુજબ
    $\cutrank{(\lambd[x^{!A}][N])Q}$ કરતાં નાનો છે.
  \item $M$નું સ્વરૂપ
    \[
    \ore{\ori{i}{}{\typeof{N}{!A_i}}}
        {\typeof{x_1}{!A_1}}{\typeof{N_1}{!C}}
        {\typeof{x_2}{!A_2}}{\typeof{N_2}{!C}},
    \]
    હોય, તો $M' \ident \Subst{N_i}{N}{\typeof{x_i}{!A_i}}$ છે.
    $M'$માંનું એક સંકોચ્ય પદ લો. કાં તો~$N_i$માં
    તેને અનુરૂપ સંકોચ્ય પદનો કટ-ક્રમ સમાન છે,
    જે ધારણા મુજબ $\cutrank{M}$ કરતાં નાનો છે;
    અથવા તેનો કટ-ક્રમ $\len{!A_i}$ છે, જે વ્યાખ્યા મુજબ
    $\cutrank{\ore{\ori{i}{}{\typeof{N}{!A_i}}}
      {\typeof{x_1}{!A_1}}{\typeof{N_1}{!C}}
      {\typeof{x_2}{!A_2}}{\typeof{N_2}{!C}}}$ કરતાં નાનો છે.
  \end{enumerate}
  \sourcecorrection{OLPG-0687-COPIED}{મૂળના બીજા અને ત્રીજા કિસ્સાની સંકોચ્ય પદોની યાદીમાં પ્રતિસ્થાપિત પદમાંથી નકલ થતા આંતરિક સંકોચ્ય ઉપપદોનો કિસ્સો સ્પષ્ટ કરેલો નથી. દાવો ચકાસવા આ વધારાનો કિસ્સો પણ જરૂરી છે; સ્રોતમાં તે અધૂરો છે.}
\end{proof}

\begin{thm}
  બધાં સાબિતી-પદો સામાન્ય સ્વરૂપમાં ઘટે છે;
  બધી !!{derivation}s સામાન્ય !!{derivation}sમાં ઘટે છે.
\end{thm}

\begin{proof}
  બીજો દાવો પહેલામાંથી અનુસરે છે. પહેલા દાવા માટે
  $m=\maxrank{M}$ પર પૂર્ણ આગમન કરીએ; અહીં $M$
  સાબિતી-પદ છે.
  \begin{enumerate}

  \item $m=0$ હોય, તો $M$ પહેલેથી સામાન્ય છે.
  \item
    નહિતર~$n$ પર આગમન કરીએ; તે~$M$માં
    કટ-ક્રમ~$m$ હોય એવાં સંકોચ્ય પદોની સંખ્યા છે.
    \begin{enumerate}
    \item $n=1$ હોય, તો એવું સંકોચ્ય પદ~$N$ પસંદ કરો
      કે સંકોચ્ય હોય એવા તેના દરેક પોતાના કરતાં
      નાના ઉપપદ~$P$ માટે $m = \cutrank{N} > \cutrank{P}$
      હોય. દરેક પદનાં ફક્ત સીમિત ઉપપદો હોય છે,
      એટલે આવું સંકોચ્ય પદ હોવું જ જોઈએ.

      $N$ના સંકોચનનું પરિણામ~$N'$ હોય.
      ઉપપાદ્યથી $\maxrank{N'} < \maxrank{N}$ છે.
      તેથી $\cutrank=m$ હોય એવાં સંકોચ્ય પદોની
      સંખ્યા~$n$ ઘટે છે. એટલે $m$ પણ ઘટે છે
      ($1$ કે વધુ જેટલું), અને~$m$ માટેની
      આગમન-ધારણા લાગુ કરી શકાય છે.
    \item આગમનના પગલા માટે $n>1$ ધારો.
      પ્રક્રિયા સમાન છે; ફક્ત~$n$ ઘટીને પણ ધન
      રહે છે, એટલે~$m$ બદલાતો નથી.
      હવે~$n$ માટેની આગમન-ધારણા લાગુ કરીએ.
    \end{enumerate}
\end{enumerate}
\sourcecorrection{OLPG-0687-CONTEXT}{મૂળની આ આગમન-દલીલ ફક્ત ઘટાડેલા ઉપપદની અંદરનો કટ-ક્રમ નિયંત્રિત કરે છે. તે ઉપપદને તેના સંકોચનના પરિણામથી બદલતાં બહારના આવરણમાં નવો સંકોચ્ય પદ ખુલ્લો પડી શકે છે, ખાસ કરીને વિકલ્પ-લોપનું પરિણામ વિધેય હોય ત્યારે. આવા નવા કટના ક્રમ અને મહત્તમ-ક્રમવાળા કટની કુલ સંખ્યા માટે મૂળ દલીલ આપતું નથી.}
\end{proof}

પ્રકારિત $\lambd$-પદોનું સામાન્યીકરણ થાય છે તેનો
અર્થ એ કે દરેક આવા પદનું \emph{સામાન્ય સ્વરૂપ છે}.
$\lambd$-પદને $\beta$-ઘટાડાથી સામાન્ય સ્વરૂપમાં
લાવવું ગણતરી કરવી છે, અને સામાન્ય સ્વરૂપ
ગણતરીનું મૂલ્ય છે એમ માનીએ, તો તેનો અર્થ
એ થાય કે પ્રકારિત $\lambd$-કલનની દરેક ગણતરી
આખરે થોભે છે અને મૂલ્ય આપે છે. વળી,
પ્રકારનું સંરક્ષણ એ મૂલ્યનો પ્રકાર યોગ્ય
હોવાની ખાતરી આપે છે. દાખલા તરીકે, $N$નો
પ્રકાર $!A \lif !B$ હોય (એટલે તે~$!A$
પ્રકારના આર્ગ્યુમેન્ટ પર કામ કરતું અને~$!B$
પ્રકારનાં મૂલ્યો આપતું વિધેય હોય) અને તેને~$!A$
પ્રકારના પદ~$M$ (ઇનપુટ) પર લાગુ કરીએ,
તો $(NM)$ પદ~$N'$ (આઉટપુટ)માં ઘટે છે,
અને તે આઉટપુટનો પ્રકાર~$!B$ હોવાની ખાતરી છે.
\sourcecorrection{OLPG-0687-STRONG}{આગળનો સબળ સામાન્યીકરણનો દાવો દરેક ઘટાડાક્રમ માટે છે, પરંતુ ઉપરની સાબિતી ફક્ત એક નિર્ધારિત પસંદગીથી સામાન્ય સ્વરૂપ મળે તે બતાવવાનો પ્રયાસ કરે છે. સબળ સામાન્યીકરણની અલગ સાબિતી મૂળમાં નથી; તેથી દરેક ગણતરી થોભવાના દાવાને ઉપરના પ્રમેયથી સાબિત ગણ્યો નથી.}

હકીકતમાં પ્રકારિત $\lambd$-કલનનાં પદો માટે
ઉપપદો પર $\beta$-ઘટાડા લગાવવાની કોઈ એક
રીતથી જ સામાન્ય સ્વરૂપ મળતું નથી; ઉપપદો
પર $\beta$-ઘટાડા લગાવવાની \emph{દરેક}
રીત આખરે સામાન્ય સ્વરૂપમાં સમાપ્ત થાય છે.
આ ગુણધર્મ \emph{સબળ સામાન્યીકરણ} કહેવાય.
તેનો અર્થ એ કે ગણતરી આખરે થોભીને મૂલ્ય
આપે તેની ખાતરી આપતી કોઈ એક રીત જ નથી,
પણ ગણતરી કરવાની \emph{દરેક} રીત
આખરે મૂલ્ય આપે છે.

ગણતરી કરવાની જુદી રીતો \emph{એક જ} મૂલ્ય
આપે છે તે કેવી રીતે જાણીએ? અત્યાર સુધી
પ્રકારના સંરક્ષણથી ફક્ત એટલું જાણીએ છીએ
કે મળતાં બધાં મૂલ્યોનો પ્રકાર સમાન છે.
પ્રકારિત $\lambd$-કલનનો બીજો અગત્યનો
ગુણધર્મ છે: તે \emph{સંગમક્ષમ}, એટલે
\emph{ચર્ચ--રોસર}, છે.
$N \red N_1$ અને $N \red N_2$ હોય,
તો એવું પદ~$N'$ છે કે $N_1 \red N'$
અને $N_2 \red N'$ થાય. યાદ કરો કે
$N \red N'$ એટલે શૂન્ય કે વધુ
$\redone$-ઘટાડા પછી $N'$ મળે છે.
$N_1$ સામાન્ય સ્વરૂપમાં હોય, તો
$N_1 \red N'$ થાય એવું એકમાત્ર પદ~$N'$
એ~$N_1$ પોતે જ છે (શૂન્ય
$\redone$-ઘટાડા વડે). તેથી $N_1$ અને
$N_2$ બંને સામાન્ય સ્વરૂપમાં હોય, તો
$N_1 = N' = N_2$ છે. એટલે~$\red$
સબળ સામાન્યીકરણ ધરાવતો હોવાથી પદને
ઘટાડવાની દરેક રીત આખરે સામાન્ય સ્વરૂપ આપે છે.
વળી~$\red$ ચર્ચ--રોસર હોવાથી કોઈ પણ
બે સામાન્ય સ્વરૂપ સમાન છે. આમ, પ્રકારિત
$\lambd$-કલનમાં ગણતરી હંમેશાં થોભે છે
અને તે કેવી રીતે આગળ વધી હોય તોપણ
મૂલ્ય, એટલે સામાન્ય સ્વરૂપ, એક જ હોય છે.

સંબંધ સંગમક્ષમ છે તે સાબિત કરવું સામાન્ય રીતે
અઘરું છે. જોકે આ નબળી શરત સહેલાઈથી
સ્થાપી શકીએ:

\begin{prop}[નબળો ચર્ચ--રોસર ગુણધર્મ]
$N \redone N_1$ અને $N \redone N_2$ હોય,
તો એવું પદ~$N'$ છે કે $N_1 \red N'$
અને $N_2 \red N'$ થાય.
\end{prop}

\begin{proof}
  $N$માં સંકોચ્ય પદ~$M_1$ અથવા~$M_2$ને
  તેના સંકોચનના પરિણામથી બદલી અનુક્રમે
  $N_1$ અને $N_2$ મળે છે. સંકોચ્ય પદો
  $M_1$ અને $M_2$ એ~$N$નાં ઉપપદો છે,
  એટલે તેઓ આંશિક રીતે પરસ્પર આવરી લેતાં નથી.
  કાં તો એક બીજાનું ઉપપદ હોય છે, અથવા
  તેઓ~$N$માં જુદાં, એકબીજાને ન આવરી લેતાં
  સ્થાનો પર આવે છે. બીજો કિસ્સો લો:
  $N$નું સ્વરૂપ $N(M_1,M_2)$ છે.
  $M_1$ને તેના સંકોચનના પરિણામ~$M_1'$થી
  બદલો, તો $N_1$ એટલે $N(M_1',M_2)$
  મળે; અને $M_2$ને તેના પરિણામ~$M_2'$થી
  બદલો, તો $N_2$ એટલે $N(M_1,M_2')$
  મળે. બંને $N(M_1',M_2')$માં ઘટે છે.
  પહેલા કિસ્સા માટે $M_2$ એ~$M_1$નું
  ઉપપદ છે એમ ધારો (બીજો કિસ્સો સમમિત છે),
  અને $M_1$ કયા પ્રકારનું સંકોચ્ય પદ છે
  તે મુજબ કિસ્સા પાડો. દાખલા તરીકે,
  $M_1 \ident \proj{1}{\pair{O_1}{O_2}}$
  હોય, તો તે $O_1$માં ઘટે છે. $M_2$
  એ~$O_1$નું ઉપપદ હોય, તો $M_2$ને
  રૂપાંતરિત કરવાથી $O_1'$ મળે છે.
  એટલે પહેલા $M_1$ અને પછી $M_2$ને
  રૂપાંતરિત કરતાં $O_1'$ મળે છે.
  પણ પહેલા $M_2$ને રૂપાંતરિત કરતાં
  $\proj{1}{\pair{O_1'}{O_2}}$ મળે છે,
  જે પણ~$O_1'$માં ઘટે છે.
  \sourcecorrection{OLMD-284}{મૂળ છેલ્લે $O_2$ કહે છે, પરંતુ પ્રથમ પ્રક્ષેપનું પરિણામ $O_1'$ જ છે, જે બીજી ઘટાડા-દિશામાં પણ મળ્યું છે.}
  \sourcecorrection{OLPG-0687-LOCAL}{મૂળની આ સાબિતી આંતરિક સંકોચન માટે ફક્ત પ્રથમ પ્રક્ષેપના રાખેલા ઘટકનો કિસ્સો આપે છે. કાઢી નાખાતા ઘટક, વિધેય-અરજી અને વિકલ્પ-લોપના અંદર-બહારના ઘટાડા તથા આગળના વિભાગના ક્રમચય-રૂપાંતરોના કિસ્સા અહીં સાબિત નથી.}
\end{proof}

\begin{lem}[ન્યૂમનનું ઉપપાદ્ય]
$\red$ સબળ સામાન્યીકરણ ધરાવતો હોય અને નબળો
ચર્ચ--રોસર ગુણધર્મ ધરાવતો હોય, તો તે સંગમક્ષમ છે.
\end{lem}

\begin{proof}
  $\red$ સબળ સામાન્યીકરણ ધરાવે છે, એટલે
  દરેક ઘટાડા-અનુક્રમ સામાન્ય સ્વરૂપ આપે છે.
  સામાન્ય સ્વરૂપ અનન્ય છે જો અને તો જ~$\red$
  સંગમક્ષમ છે. તેથી ધારો કે કોઈ પદ~$N$ને
  બે જુદાં સામાન્ય સ્વરૂપો $N'$ અને $N''$
  છે, જે આ ઘટાડા-અનુક્રમોથી મળે છે:
  \begin{align*}
  N & = N_0' \redone N_1' \redone \dots \redone N_k' = N' \text{ અને}\\
  N & = N_0'' \redone N_1'' \redone \dots \redone N_l'' = N''
  \end{align*}
  (ઘટાડા-અનુક્રમની લંબાઈ $>0$ હોવી જ જોઈએ;
  નહિતર $N$ પહેલેથી સામાન્ય સ્વરૂપમાં હોત.)
  \sourcecorrection{OLMD-285}{મૂળ બંને અનુક્રમમાં શરૂઆતને અનુક્રમણિકાંક 1 આપે છે, પરંતુ આગળની દલીલ 1વાળાં પદોને પ્રથમ ઘટાડા પછીનાં પદો માને છે. શરૂઆતને 0 અને પ્રથમ ઉત્તરગામીને 1 આપ્યાં છે.}

  $N_1' = N_1''$ હોય, તો તેને બે જુદાં
  સામાન્ય સ્વરૂપો ($N'$ અને $N''$) છે.
  $N_1' \neq N_1''$ હોય, તો નબળા
  ચર્ચ--રોસર ગુણધર્મથી એવું $N'''$ મળે છે
  કે $N_1' \red N'''$ અને $N_1'' \red N'''$
  થાય. $N' \neq N''$ હોવાથી,
  કાં તો $N''' \neq N'$ અથવા
  $N''' \neq N''$ છે. પરિણામે,
  $N_1'$ અથવા $N_1''$ને બે જુદાં
  સામાન્ય સ્વરૂપો છે. આપણે બતાવ્યું કે
  $N$ને બે જુદાં સામાન્ય સ્વરૂપો હોય,
  તો એવું $N_1$ છે કે $N \redone N_1$
  થાય અને $N_1$ને પણ બે જુદાં સામાન્ય
  સ્વરૂપો હોય. આ રીતે આગળ વધીને
  $N \redone N_1 \redone N_2 \redone \dots$
  મેળવી શકાય. પરંતુ~$\redone$ સબળ
  સામાન્યીકરણ ધરાવે છે, એટલે આ અશક્ય છે.
  \sourcecorrection{OLPG-0687-NEWMAN}{મૂળની આ દલીલમાં સામાન્ય સંયુક્ત ઉત્તરગામી $N'''$ પોતે સામાન્ય સ્વરૂપ હોવાની ખાતરી નથી. બે જુદાં સામાન્ય સ્વરૂપોનો દાવો આગળ વધારવા તેના સામાન્ય સ્વરૂપ સુધી ઘટાડાનો પગલો જરૂરી છે; મૂળ તે પગલો આપતું નથી.}
\end{proof}


\end{document}
